\chapter{Introduction}

\section{Background}

Large language models (LLMs) have made significant progress in the domain of natural language processing (NLP). LLMs have demonstrated exceptional competence in generating text, summarizing texts, and answering questions. However, LLMs are heavily dependent on parametric knowledge learned through training, thereby limiting their ability to retrieve relevant information and resulting in problems like hallucinations and lack of factual consistency.

In order to address these challenges, the paradigm of retrieval augmented generation (RAG) has been proposed. The RAG model involves the combination of LLMs with external knowledge sources. In this approach, relevant documents are first retrieved using a vector database and then used as context for generation.

\section{Problem Statement}

Even though RAGs possess various strengths, the efficacy of the model is significantly influenced by the accuracy of its retrieval mechanism. One of the primary shortcomings arises when the system processes syntactically ambiguous sentences. Traditional methods rely on keyword searching and vector similarity, which fail to account for deeper linguistic dependencies in the language. As a result, the retrieval process often returns incorrect documents, leading to misguiding output.

The issue becomes increasingly prominent in situations where there are instances of prepositional phrase attachment and reduced relative clauses. The classical approach to natural language processing suffers from a local attachment bias. This implies that decisions are made based on proximity rather than global sentence structures.

\section{Motivation}

The incapacity of classical RAG models to disambiguate syntactic ambiguity highlights the need for more powerful computational models. Quantum computing appears to be an appealing choice due to its ability to describe and compute data within Hilbert space. Through concepts like superposition and entanglement, quantum-based algorithms can model intricate correlations that are difficult to handle using classical models.

This research is guided by the premise that quantum models may enhance the performance of classical retrieval-based models, especially when dealing with non-linear structures in language. Instead of replacing conventional models, the aim is to develop a combination framework utilizing the best attributes of both worlds.

\section{Proposed Approach}

QRAG is introduced in this paper as the proposed Quantum RAG approach, which combines traditional RAG systems with Quantum Natural Language Processing (QNLP). The system includes the Ambiguity Controller, which checks for structurally complex questions and sends them to the quantum pipeline for processing.

The quantum pipeline uses the Variational Quantum Circuit (VQC) for the encoding of grammatical dependencies as well as syntactic disambiguation. On the other hand, simpler questions will be processed using classical methods efficiently.

\newpage

\section{Contributions}

The key contributions of this work are as follows:

\begin{itemize}
    \item Identification of syntactic ambiguity as a critical limitation in classical RAG systems.
    \item Design of a hybrid Quantum-Classical framework for improved semantic understanding.
    \item Integration of Variational Quantum Classifiers for syntactic disambiguation.
    \item Development of an Ambiguity Controller for efficient query routing.
    \item Empirical validation demonstrating improved accuracy on ambiguous query datasets.
\end{itemize}

\section{Organization of the Report}

This report will be divided as follows. The upcoming chapters will explain the theories behind RAG and quantum computing before describing the approach used by QRAG in detail. Implementation issues and a demonstration of the approach will follow, after which a discussion on how to move forward and conduct further research will conclude the report.